\documentclass{article}

% Adds the url command
\usepackage{hyperref}
% To remove paragraph indentation
\usepackage{parskip}
% Adds line cut to urls
\usepackage{xurl}

% Update these:
\title{Lab 6 Report}
\date{Due: 03/10/2026 @ 11:00 AM}
\author{Ryan Wilson}
\begin{document}
\maketitle
\newpage


\section{Introduction}
The purpose of this lab report is to determine the time complexity of four numerical integration methods: the Midpoint Rule, Simpson's 1/3 Rule, Simpson's 3/8 Rule, and the Gauss Quadrature. To find the Big O complexity, analyzing the number of operations required for each method and then finding the growth rate is necessary.

\section{Midpoint Rule}
The midpoint of the function is calculated as: $f(x_{mp}) = f((a + b)/2)$\\\\
The function is calculated one time. Since the method has a fixed number of calculations regardless of the size of the function, the number of lines of code executed is constant. Therefore, the Big O complexity is $\mathcal{O}(1)$.\\

\section{Simpson's 1/3 Rule}
The Simpson's 1/3 Rule calculates the function at 3 different points: the left bound, the midpoint, and the right bound. Simpson's 1/3 Rule is calculated as: $f(x_{s, 1/3}) = f(a) + 4f((a + b)/2) + f(b)$\\\\
The function is calculated one time. Since the method has a fixed number of calculations regardless of the size of the function, the number of lines of code executed is constant. Therefore, the Big O complexity is $\mathcal{O}(1)$.\\

\section{Simpson's 3/8 Rule}
The Simpson's 3/8 Rule calculates the function at 4 different points: the left bound, two internal points that are evenly spaced, and the right bound. Simpson's 3/8 Rule is calculated as: $f(x_{s, 3/8}) = f(a) + 3f((2a + b)/3) + 3f((a + 2b)/3) + f(b)$.\\\\
The function is calculated one time. Since the method has a fixed number of calculations regardless of the size of the function, the number of lines of code executed is constant. Therefore, the Big O complexity is $\mathcal{O}(1)$.\\

\section{Gauss Quadrature}
The Gauss Quadrature, or Gauss Quad, uses a loop to iterate through a selected number n of quad points. The code calculates the function for each value of n, staying between the bounds of 0 to 10. Gauss Quad is calculated as: $x_{gq} = f(c + mt_i)$ where $c = (a + b)/2$, $m = (b - a)/2$, and t is the abcissa, or horizontal value.\\\\
The function uses a loop that runs from i=1 to n. Each loop iteration calculates a constant number of operations, meaning the total number of lines of code executed is proportional to n. Therefore, the Big O complexity is $\mathcal{O}(n)$.\\

\section{Sources}
\begin{itemize}
\item AERE361 Lab Manual 6
\item TAs
\item \url{https://pomax.github.io/bezierinfo/legendre-gauss.html}
\end{itemize}

\end{document}
